\documentclass[11pt,a4paper]{article}
\usepackage[T1]{fontenc}
\usepackage{lmodern,amsmath,amssymb,amsthm,mathtools}
\usepackage[margin=0.82in]{geometry}
\usepackage{microtype,booktabs,longtable,array,enumitem,fancyhdr,xurl}
\usepackage{xcolor}
\usepackage[colorlinks=true,linkcolor=blue!50!black,urlcolor=blue!50!black,citecolor=blue!50!black]{hyperref}
\newtheorem{theorem}{Theorem}[section]
\newtheorem{lemma}[theorem]{Lemma}
\newtheorem{proposition}[theorem]{Proposition}
\newtheorem{corollary}[theorem]{Corollary}
\theoremstyle{remark}\newtheorem{remark}[theorem]{Remark}
\newcommand{\Z}{\mathbb Z}
\newcommand{\N}{\mathbb N}
\newcommand{\Con}{\operatorname{Con}}
\newcommand{\ZFC}{\mathsf{ZFC}}
\newcommand{\ZF}{\mathsf{ZF}}
\newcommand{\Ccal}{\mathcal C_{222}}
\newcommand{\vp}{v_2}
\newcommand{\htp}{\operatorname{ht}}
\newcommand{\pair}[2]{\langle #1,#2\rangle}
\setlength{\parskip}{0.35em}
\setlength{\parindent}{0pt}
\setlength{\headheight}{14pt}
\pagestyle{fancy}\fancyhf{}
\fancyhead[L]{\small ProveIt: Diophantine consistency}
\fancyhead[R]{\small Pell growth and an affine radix}
\fancyfoot[C]{\thepage}
\hypersetup{pdftitle={Pell growth from divisibility: an explicit 210-variable quartic for Con(ZFC)},pdfauthor={Research prepared for Vladimir Reshetnikov}}
\allowdisplaybreaks[1]
\begin{document}
\hypersetup{pageanchor=false}
\begin{titlepage}
\vspace*{1.4cm}
{\Large\scshape ProveIt research report\par}
\vspace{1.0cm}
{\huge\bfseries Pell growth from divisibility\par}
\vspace{0.35cm}
{\LARGE An affine CRT radix and an explicit\par 210-variable quartic for $\Con(\ZFC)$\par}
\vspace{0.8cm}
{\large Prepared for Vladimir Reshetnikov\par 25 October 2026\par}
\vspace{1.0cm}
\begin{abstract}
We construct a literal quartic polynomial in 210 unrestricted integer variables
whose non-solvability is equivalent, provably in ZFC, to the consistency of ZFC,
by composing new arithmetic reductions with the repository's identified
numerical-calculus calibration. The polynomial is a sum of 150 squares of
quadratic polynomials. Its collected expansion has 5,282 nonzero monomials and
coefficient height $2^{40}$. All coefficients and all 150 residuals are supplied.

The main change replaces an exact exponential-growth request by three
quadratic equations forcing the existing growth coordinate to be a sufficiently
large Pell ordinate. A normalized Pell permutation modulo powers of two makes
this replacement compatible with the odd binomial selector required for CRT.
The existing dominant CRT index also provides an affine radix base, removing
another coordinate. A delayed sign argument is essential: the new growth
ordinate is initially signed. The construction improves the incoming
213-variable quartic in variable count, number of residuals, support and
coefficient height. We also provide a 214-variable companion with much smaller
coefficients, and an independent 11-positive-variable, degree-174 companion.

These are explicit constructions and ordinary mathematical proofs, with exact
finite coefficient and reconstruction checks. No global minimum, historical
priority, fresh end-to-end Lean/Rocq build, or proof of consistency is claimed.
\end{abstract}
\vfill
\textbf{Principal equation:}
\[
P_{210}(x_0,\ldots,x_{209})=\sum_{j=0}^{149}R_j(x_0,\ldots,x_{209})^2=0,
\qquad \boldsymbol x\in\Z^{210}.
\]
\textbf{Complete literal residuals:} Appendix~\ref{app:literal}.\par
\textbf{Executable specification:} \texttt{equations/principal210/equation.json}.\par
\textbf{Offline verification:} \texttt{python3 -B verify.py}.
\end{titlepage}
\hypersetup{pageanchor=true}
\pagenumbering{roman}
\tableofcontents
\clearpage
\pagenumbering{arabic}

\section{The result and its exact scope}

The word ``simplest'' involves several competing measurements. A small number
of variables can require a large degree or large fixed coefficients; a
quartic with small coefficients can require many variables. This report
keeps these measurements separate and makes a concrete improvement in the
direct quartic family.

\begin{theorem}[Explicit principal equation]\label{thm:main}
Let $R_0,\ldots,R_{149}\in\Z[x_0,\ldots,x_{209}]$ be the literal quadratic
polynomials in Appendix~\ref{app:literal}, and put
$P_{210}=\sum_{j=0}^{149}R_j^2$. Then $P_{210}$ has exact degree four,
5,282 collected nonzero monomials and coefficient height
$2^{40}=1,099,511,627,776$. Composing the constructions proved here with the
explicitly identified inherited arithmetic and logical calibration gives
\[
\ZFC\vdash
\left(\neg\exists\boldsymbol x\in\Z^{210}\;P_{210}(\boldsymbol x)=0\right)
\ \longleftrightarrow\ \Con(\ZFC).
\]
All 210 coordinates range independently over the unrestricted integers.
\end{theorem}

Consistency here means the absence of a finite syntactic refutation in the
effective ZFC presentation used by the inherited calibration. The proof is
internal to ZFC: proof codes and integer witnesses range over the natural
numbers and integers of the theory. It does not introduce a restriction to
externally standard codes. We prove the equivalence without assuming either
consistency or inconsistency, and we do not prove the polynomial has no zero.

\subsection{Measured comparisons}

The repository snapshot inspected for this report is
\texttt{ef0c534c1ae1e312c61a2a9a5a3f18f5516fd42b}. Its incoming archive named
\texttt{ConZFC\_CRT\_Index\_Anchor\_215\_2026-10-25.zip} already contains
213- and 214-variable companions. Thus 213, rather than 215, is the relevant
existing variable minimum for this comparison. The frozen predecessor is
included unchanged under \texttt{inherited/anchor/}.

\begin{center}\small
\begin{tabular}{lrrrrr}
\toprule
Equation & Variables & Rows & Degree & Monomials & Height\\
\midrule
Incoming balanced & 215 & 154 & 4 & 5,660 & 91,679\\
Incoming companion & 214 & 153 & 4 & 5,486 & $2^{40}+76$\\
Incoming count minimum & 213 & 152 & 4 & 5,402 & $2.41785\ldots\times10^{24}$\\
New balanced companion & 214 & 153 & 4 & 5,699 & 91,679\\
\textbf{New principal} & \textbf{210} & \textbf{150} & \textbf{4} & \textbf{5,282} & $\boldsymbol{2^{40}}$\\
\bottomrule
\end{tabular}
\end{center}

The exact height of the incoming 213-variable equation is
$2,417,851,639,242,452,488,945,685$. The new principal improves all five
reported size quantities except degree, which is unchanged, relative to that
equation; it also improves the incoming 214-variable companion in all four
non-degree measurements. The balanced 215-variable predecessor still has
much smaller coefficients. The new balanced 214-variable companion retains
its height while trading 39 additional monomials for one fewer coordinate.
It applies only the affine-radix substitution to the original balanced source,
retaining that source's original growth certificate. The three-row growth
replacement belongs to the principal 210-variable construction.

The maximum coefficient in a principal quadratic residual is $2^{20}$.
The monomial $x_0^2x_{209}^2$ has collected coefficient $2^{40}$, proving that
the reported quartic degree and height are attained. These figures are computed
after collecting equal monomials, not from an unexpanded circuit estimate.

\subsection{What is new and what is inherited}

The use of Pell sequences for Diophantine growth is classical; exponential
growth was central to Matiyasevich's work \cite{matiyasevich}. The contribution
here is a compatible replacement inside a specific fully instantiated
consistency construction. It consists of a three-row growth certificate,
an affine radix recovered from an existing divisibility, a sign recovery
argument, and a cofinal binary-residue construction needed for completeness.
The existing common-Hessian and affine projections are reused, with new
certificates. Elementary integer coordinate and residual-row operations
reduce the final support.

The principal construction works directly with the calibrated numerical
calculus $\Ccal$. It does not simulate K215 and does not instantiate Jones'
universal polynomial. The separate compact companion in
Section~\ref{sec:compact} retains its own native/compiler calibration.

\section{The inherited arithmetic interface}\label{sec:interface}

We describe the semantic coordinates before the exact projections. They are
not all separate coordinates after projection; the final equation is defined
literally in the appendix.

The positive horizon has the cofinal form
\begin{equation}\label{eq:horizon}
t=16+4z^2.
\end{equation}
The coordinate $F$ provides formula weights $1+jF$. The quantity
$B=\binom{F+t}{t}$ will be recovered from a radix digit, and $N=B^x$ is the
generalized CRT modulus. The beta spacing is $y=F(2v+1)$ and the beta code
$x$ satisfies
\begin{equation}\label{eq:betaend}
x=F+1+(y+1)r_0,
\qquad x=1+[1+(t+1)y]r_t.
\end{equation}
Eight independent CRT representatives encode the four local parameters, two
pending states, and two beta quotients. The ninth representative $C$ encodes
the row index. It is explicitly even and satisfies
\begin{equation}\label{eq:selector}
B\eta=C+F+1.
\end{equation}

The shared Pell machinery has its own large center $A_*$ and a square frame
containing every actual requested base, exponent and output. Its inherited
signed index theorem first authenticates the principal Pell pairs at their
nominal, initially signed indices. Its distinguished anchor is the pair for
$C$, and its cap proves
\begin{equation}\label{eq:cap}
C\ne0,\qquad |k|<|C|\quad\hbox{for every other retained nominal index }k.
\end{equation}
After the new growth pair is removed there are 13 retained principal pairs
and 15 exact-power requests. These include the requests at base $q+1$ with
indices $F+t$, all eight ordinary representatives, and $C$; the requests at
base $q$ with indices $t$, $y+1$, and $x$; and the two singleton requests
with exponent $x$ and bases $B$ and the radix base.

For any integer base $b$ with $|b|\ge2$, the inherited signed power theorem
forces a nonnegative requested exponent $k$ and the exact output $b^k$.
This theorem is used only after the main pairs have been oriented. The frame
and anchor theorem does not require the front coordinates to have recovered
their arithmetic signs. The new independent Pell witnesses will not be
members of this frame; their existing output coordinate $F$ remains a frame
atom.

The exact finite changes to this interface are checked by independently
reconstructing all 34 retained square-frame atoms, all 13 principal norms and
multiplication wires, the whole orientation sum and the anchor cap.

\subsection{The literal local syntax}

The inherited square-shell constructors are
\begin{align*}
I(a,b)&=2((a+b)^2+a),&M(a,b)&=I(a,b)+1,\\
A(a,b)&=2((a+b)^2+3a+4b+3),&E(a,b)&=A(a,b)+1.
\end{align*}
They stand for implication, membership, universal quantification and equality.
The numerical truth and falsity codes are 0 and 1. The incoming recursive
recoding theorem transports the original numerical calculus to these
constructors, including variable labels and the guarded seed schema. This
recoding and its finite constructor expressions are retained unchanged.

The source-bound constructor audit gives maximum depth 19 for a local formula
label. The nine axiom terms have depths
\[
4,3,3,3,3,5,4,3,19.
\]
Every leaf parameter is strictly below $t$. The following sharper majorant
uses the same expression trees; it does not change the calculus.

\begin{lemma}[A capacity bound for constructor trees]\label{lem:capacity}
Put $d=2^{19}$ and $L=8^{d-1}t^d$. Every local formula label is strictly
less than $L$.
\end{lemma}
\begin{proof}
For $0\le u,v\le m$, the largest constructor is bounded by $E(m,m)$, and
\[
E(m,m)+1=8m^2+14m+8\le8(m+1)^2.
\]
A depth-zero leaf has value plus one at most $t$. The recurrence
$H_0=t$, $H_{r+1}=8H_r^2$ therefore bounds value plus one at depth $r$.
Its explicit solution is $H_r=8^{2^r-1}t^{2^r}$. At depth 19 it equals
$L$. Unequal child depths are handled by monotonicity. Thus a label has value
at most $L-1$, as asserted.
\end{proof}

\section{Growth from three quadratic equations}\label{sec:growth}

The old exact request $F=t^{2^{20}t}$ is stronger than the proof extraction
needs. The required interface is a sufficiently large even nonzero $F$,
eventually positive, with a cofinal completeness construction. Introduce
three integer coordinates $X,Z,u$ and impose
\begin{equation}\label{eq:growth}
\boxed{\quad Z=tF,\qquad X^2-16Z^2-8ZF=1,\qquad
F=2^{19}t(2u+1).\quad}
\end{equation}
These are quadratic equations. They replace five coordinates of the old
growth request: its principal Pell triple, its modulus and its quotient.
The obsolete growth index is removed from the frame. The old principal pair's
norm, multiplication wire, orientation term and cap term are removed exactly
as specified by the source certificate; $F$ remains a frame atom.

\subsection{Pell recurrence facts}

For $a\ge2$ define $X_n(a),Y_n(a)$ by
\begin{align*}
X_0&=1,&Y_0&=0,\\
X_{n+1}&=aX_n+(a^2-1)Y_n,&Y_{n+1}&=X_n+aY_n.
\end{align*}
The least positive Pell unit at discriminant $a^2-1$ is
$a+\sqrt{a^2-1}$: the ordinate of a nontrivial positive integral solution
is at least one, and $(a,1)$ is a solution. The standard Pell classification
therefore gives, for every integer solution with $Y\ne0$,
$|X|=X_n(a)$ and $|Y|=Y_n(a)$ at a unique $n\ge1$.

We need three elementary consequences:
\begin{align}
Y_n(a)&\equiv n\pmod{a-1},\label{eq:indexcong}\\
\vp(Y_n(a))&=\vp(n) &&(a\hbox{ odd},\ n\ge1),\label{eq:valuation}\\
Y_n(a)&\ge(2a-1)^{n-1} &&(n\ge1).\label{eq:pellgrowth}
\end{align}
The first follows by specializing the integer polynomial recurrence at
$a=1$. For the second, all $X_n(a)$ are odd and $Y_n(a)$ has the parity of
$n$ when $a$ is odd; the duplication identity
$Y_{2m}=2X_mY_m$ adds exactly one factor of two at each doubling. For the
third, $Y_n$ is increasing and
$Y_{n+1}=2aY_n-Y_{n-1}\ge(2a-1)Y_n$.

\begin{theorem}[Sign-free growth]\label{thm:growth}
Every integer solution of \eqref{eq:horizon} and \eqref{eq:growth} has even
nonzero $F$ and
\[
|F|\ge L^t,\qquad L=8^{2^{19}-1}t^{2^{19}}.
\]
In particular $|F|\ge t^t>t$.
\end{theorem}
\begin{proof}
The last equation makes $F$ nonzero. Eliminating $Z$ gives
\[
X^2-((4t+1)^2-1)F^2=1.
\]
Pell classification at the odd center $a=4t+1$ gives
$|F|=Y_n(a)$ for $n\ge1$. Put $d=2^{19}$. Since $4t\mid F$,
\eqref{eq:indexcong} gives $4t\mid n$. Moreover,
\[
\vp(n)=\vp(F)=19+\vp(t).
\]
Consequently the positive integer $n/(4t)$ has exact valuation 17, and
$n\ge dt$. Formula \eqref{eq:pellgrowth} gives
\[
|F|\ge(8t+1)^{n-1}\ge(8t)^{dt-1}\ge
\bigl(8^{d-1}t^d\bigr)^t.
\]
The last ratio is exactly $8^{t-1}/t\ge1$. This proves the result with no
assumption about the sign of $F$.
\end{proof}

The flexibility of this certificate is useful. It does not insist that $F$
be one specified power. Section~\ref{sec:cofinal} will choose among its Pell
ordinates to meet an additional binary-residue condition. Without that
choice, merely proving a growth lower bound would not establish completeness.

\section{The dominant CRT index supplies the radix}\label{sec:radix}

Replace the old base $B[2+4(C+\sum A_i)]$ by
\begin{equation}\label{eq:newradix}
D=2(C+F+1),\qquad q=D^x.
\end{equation}
The power in this display is still represented by the retained polynomial
certificate. It is not an operation in the final polynomial. The definition
of $D$ is affine, so its old multiplication coordinate and defining residual
can be eliminated without raising degree. Crucially, the actual affine
expression is substituted into the main square frame as well.

Equation \eqref{eq:selector} gives $B\mid D$ directly. Once the signs have
been established, the existing anchor supplies $D>2C$ and bounds every
needed binomial upper argument by $C$. These two roles replace the old
explicit factor $B$ and the old sum-majorant.

\subsection{Recovering the signs in the correct order}\label{sec:signs}

The main shared-Pell frame and $C$-anchor theorem apply first, while the front
coordinates are signed. Theorem~\ref{thm:growth} independently gives even
nonzero $F$ with $|F|>t$. Because $C$ is even, $C+F+1$ is odd and nonzero;
thus $|D|\ge2$ and the request $q=D^x$ forces $x\ge0$.

Since $y=F(2v+1)$, $|y|\ge|F|$. If $x=0$, the terminal endpoint in
\eqref{eq:betaend} would make $1+(t+1)y$ divide $-1$. This is impossible:
\[
|1+(t+1)y|\ge(t+1)|y|-1>1.
\]
Hence $x\ge1$ and $|q|\ge2$. The retained power at base $q$ and exponent
$y+1$ now forces $y+1\ge0$. Negative $y$ would be at most $-|F|<-1$.
Therefore $y\ge|F|>0$.

The value $x=1$ needs a separate exclusion. The first endpoint would give
$-F=(y+1)r_0$, whereas $0<|F|\le y<y+1$. This is impossible. Thus
$x\ge2$, and the terminal endpoint gives
\begin{equation}\label{eq:xbound}
x\ge2+(t+1)y.
\end{equation}
It follows that $|q|\ge2^x\ge4$ and $|q+1|\ge3$. The base-$(q+1)$
requests now force all actual CRT representatives and $F+t$ to be
nonnegative. If $F$ were negative, then $F+t\le t-|F|<0$; hence $F>0$.
The already nonzero $C$ is positive as well. Thus $D>0$ and $q>0$.

Finally, if $r_0\le-1$, the first endpoint gives $x\le F-y\le0$.
Therefore $r_0\ge0$ and $x\ge F+1$. We have recovered every required sign
without adding a sign witness. In particular, the old assumption of a positive
growth output has been replaced by a proved consequence of retained equations.

\subsection{One binomial bound for the whole front}

By \eqref{eq:cap}, the initial upper argument $F+t$ and all ordinary
representatives are below $C$; the anchor argument itself is $C$.
The lower cutoffs $t,y+1,x$ are at most $x$. Also
\[
D=2(C+F+1)>2C,\qquad q=D^x>(2C)^x.
\]
For every binomial coefficient in any required low digit, including $B$,
\begin{equation}\label{eq:uniformbin}
0\le\binom{a}{j}\le C^x<q/2^x
\quad(0\le a\le C,\ 0\le j\le k\le x).
\end{equation}
Coefficients with $j>a$ are zero; $j=0$ is included since $C>t\ge16$.
As $x\ge F+1$, this implies $C^x/(q-1)<1/32$.

For clarity, the retained remainder certificate uses $M=q^x$, denominators
$q^k$, integer ratios $a_k=M/q^k$, and normalized remainders $\rho_k$.
Its exact division equations imply $a_k\mid\rho_k$. Its ternary norm budget
\[
M^2=1+4\sum_k\rho_k^2+r^2+s^2+w(w+1)
\]
gives $|\rho_k|<M/2$, hence the actual integer remainder
$r_k=\rho_k/a_k$ satisfies $|r_k|<q^k/2$. The genuine binomial low tail is
less than $q^k/32$ by \eqref{eq:uniformbin}. Congruence modulo $q^k$
therefore forces the two remainders to agree exactly. This is an equality
argument, not an approximation.

The first digit consequently gives $B\equiv\binom{F+t}{t}\pmod q$.
Equation \eqref{eq:selector}, with nonzero integer $\eta$, also gives
\[
|B|\le C+F+1=D/2<q/2.
\]
The true binomial coefficient is positive and below $q/32$, so
\begin{equation}\label{eq:Bexact}
B=\binom{F+t}{t}>0.
\end{equation}
It is at least two. The retained request authenticates $N=B^x$, and
\[
q=D^x=(2B\eta)^x=N(2\eta)^x.
\]
Thus $N\mid q$, the exact modulus relation required by the generalized
range compiler. The same tail argument authenticates all remaining range
digits. The inherited local arithmetic bound below $2^x$ continues to hold:
it needs $F\ge t^t$, $t\ge16$ and \eqref{eq:xbound}, all now proved.

\section{A cofinal choice with an odd binomial selector}\label{sec:cofinal}

Soundness alone is not enough. The old exact power supplied high powers of
two automatically. The new ordinate has a prescribed finite valuation in our
completeness construction, so we must choose its other binary digits.

\begin{lemma}[Pell permutation modulo powers of two]\label{lem:permutation}
If $a\equiv1\pmod4$, the map $s\mapsto Y_s(a)$ induces a permutation of
$\Z/2^r\Z$ for every $r\ge1$. Each odd residue has arbitrarily large
positive odd indices realizing it.
\end{lemma}
\begin{proof}
All $X_s(a)\equiv1\pmod4$ and $Y_s(a)\equiv s\pmod2$. Put $h=2^r$.
Equation \eqref{eq:valuation} gives $\vp(Y_h)=r$. The norm and
$8\mid a^2-1$ imply $\vp(X_h^2-1)\ge2r+3$. Since $X_h+1$ has valuation
one, $X_h\equiv1\pmod{2^{2r+2}}$. The addition identity then gives
\[
Y_{s+h}-Y_s=(X_h-1)Y_s+X_sY_h\equiv2^r\pmod{2^{r+1}}.
\]
The two index lifts $s$ and $s+2^r$ produce the two distinct output lifts.
Induction from modulo two proves bijectivity and periodicity. Positive lifts
in the resulting residue class are cofinal and preserve odd parity.
\end{proof}

\begin{theorem}[Compatible cofinal growth witnesses]\label{thm:cofinal}
There are cofinally many horizons of the form \eqref{eq:horizon} for which
\eqref{eq:growth} has arbitrarily large positive solutions $F$ such that
$\binom{F+t}{t}$ is odd.
\end{theorem}
\begin{proof}
Choose the free horizon parameter $z$ divisible by $2^{12}$. These horizons
remain cofinal and satisfy $t\equiv16\pmod{2^{26}}$. Thus $\vp(t)=4$ and
the binary digit of $t$ at position 23 is zero.

Put $d=2^{19}$, $n_0=dt$, $y_0=Y_{n_0}(4t+1)$ and
$a_0=X_{n_0}(4t+1)$. The congruence modulo $4t$ and
\eqref{eq:valuation} show that $b_0=y_0/(dt)$ is an odd integer.
Moreover $a_0\equiv1\pmod4$, since $n_0$ is even and
$X_{2m}=2X_m^2-1$ with $X_m$ odd.

Choose $\ell>26$ with $2^\ell>t$, and set $r=\ell-23$. The odd factors
$t/16$ and $b_0$ are units modulo $2^r$. By Lemma~\ref{lem:permutation},
choose a positive odd $s$, as large as desired, such that
\[
(t/16)b_0Y_s(a_0)\equiv1\pmod{2^r}.
\]
Set $n=n_0s$ and
\[
F=Y_n(4t+1)=y_0Y_s(a_0),\quad X=X_n(4t+1),\quad Z=tF,
\quad u=\frac{F/(dt)-1}{2}.
\]
Composition of Pell units proves the displayed identity for $F$, and the odd
quotient makes $u$ integral. These values satisfy all three equations and
\begin{equation}\label{eq:binarytarget}
F\equiv2^{23}\pmod{2^\ell}.
\end{equation}
Since $t<2^\ell$ and its bit 23 is zero, the binary supports of $F$ and $t$
are disjoint. There are no carries in their addition. Writing $s_2(m)$ for
the binary digit sum, the elementary factorial identity gives
\[
\vp\binom{F+t}{t}=s_2(F)+s_2(t)-s_2(F+t)=0.
\]
Thus the binomial is odd. Cofinal positive lifts of $s$ give arbitrarily
large $F$ at each chosen horizon.
\end{proof}

The congruence proof is finite arithmetic at an arbitrarily chosen modulus;
no analytic inverse theorem or computation of an infinite 2-adic object is
needed. The executable modular examples in the archive check this construction
at the actual $d=2^{19}$, but the theorem follows from the recurrence proof
for every modulus.

\subsection{The beta moduli need only prime support}

The new construction has $\vp(F)=23$ at the selected horizons. We must not
reuse the old sufficient condition $t!\mid F(2v+1)$ for arbitrary large $t$.
Fortunately the beta CRT needs only every prime at most $t$ to divide $y$.

Choose a sufficiently large positive odd $w$ divisible by all odd primes at
most $t$, and put $y=Fw$. The integer $y$ is divisible by every prime at
most $t$. If a prime $p$ divided both $1+iy$ and $1+jy$ for
$1\le i<j\le t+1$, it would be coprime to $y$ and divide $j-i$.
Then $p\le t$, contradicting $p\mid y$. Thus the beta moduli are pairwise
coprime. The odd multiplier can also make $y$ exceed every finite pending
state and $F$, and it has exactly the required form $2v+1$.

This replacement of an unnecessarily strong factorial divisibility is part
of the completeness proof. Keeping the old assertion would leave a real gap.

\section{From a zero to a refutation, and back}\label{sec:calibration}

\subsection{Sharp separation of the weighted labels}

\begin{lemma}[Finite-product separation]\label{lem:separation}
Let $L\ge t\ge1$ and $F\ge L^t$. If two multisets of labels in $[1,L]$
have at most $t$ and $2t+1$ elements respectively and their products of
$1+jF$ are equal, then the multisets are equal.
\end{lemma}
\begin{proof}
If the factor counts differ, their smaller count $r$ is at most $t$.
The larger product is at least $(F+1)^{r+1}$; the smaller is at most
$(1+LF)^r\le L^r(F+1)^r$. Equality would imply
$F+1\le L^r\le F$, a contradiction. The counts therefore agree at some
$m\le t$.

Replace $F$ by an indeterminate $T$. Each coefficient of either product
is a positive integer at most
$\binom mk L^k\le m^{m-k}L^k\le L^m\le F$.
For $m=0$ both products are one. For $m>0$, corresponding coefficients of
the difference have absolute value strictly less than $F$, including when
one of them equals $F$: the other is positive. A nonzero difference cannot
vanish at $F$, since division by its lowest occurring power of $F$ and
reduction modulo $F$ would make a nonzero coefficient of magnitude below
$F$ divisible by $F$. Thus the polynomials agree, and unique factorization
in $\mathbb Q[T]$ recovers the label multiplicities.
\end{proof}

The decoded local Horn rows have weighted updates
\[
S_{i+1}(1+Fv_i)=S_i(1+Fp_i)(1+Fq_i),
\]
with missing premises filled by the truth label zero. The endpoint $y=F$
would give an initial beta residue zero and terminal residue one. Positive
weights propagate a zero initial state to zero at every row, so this case
is impossible. Hence $y>F$, the actual initial state is $F+1$ and the final
state is one. All states are positive. Multiplying the updates and cancelling
states gives the equality between the conclusion multiset and the premise
multiset plus the target label one. Theorem~\ref{thm:growth} and
Lemma~\ref{lem:capacity} give precisely the hypotheses of
Lemma~\ref{lem:separation}.

Cycles in a selected multiset of Horn rules do not cause a gap. Let $U$ be
the labels not derivable from the available axioms using those rules. Every
selected rule whose conclusion belongs to $U$ has at least one premise in
$U$. If the target one belonged to $U$, multiplicity balance would say
that the number of conclusions in $U$ is one more than the number of
premises in $U$, contradicting the preceding inequality. Thus there is a
finite $\Ccal$ refutation.

\subsection{Soundness of the full arithmetic source}

From an arbitrary integer zero, first classify the standalone Pell pair to
obtain even, nonzero large $F$. Independently use the retained main frame,
anchor and orientation theorem to recover the canonical signed indices.
Section~\ref{sec:signs} then proves the front signs. The uniform digit
argument recovers $B$, $N$, and all range binomials. The inherited generalized
binomial-CRT theorem applies because $F\ge L^t\ge t^t$ and hence
$F\ge\operatorname{lcm}(1,\ldots,t)^2$. Its two required conclusions are
retained: each row has a full prime-power selector exceeding $t$, and every
full prime power of $B$ is assigned to some row. The local arithmetic bound
below $2^x$ converts those congruences into actual local equations.

This decodes the beta states, valid local rules, and their strict parameter
ranges. The preceding product and Horn-flow argument yields a finite
$\Ccal$ refutation. All semantic source changes have now been accounted for;
the rest of the source decoding is the identified inherited theorem.

\subsection{Constructive completeness and the order of choices}

Start with a finite $\Ccal$ refutation. The choices must be made in the
following order; this prevents a circular size or congruence argument.
\begin{enumerate}[leftmargin=*,itemsep=0.35em]
\item Choose a cofinal horizon as in Theorem~\ref{thm:cofinal}, exceeding
the refutation length and all finite local parameters; pad by truth steps.
Choose positive $F$ and its standalone Pell witnesses with the odd binomial
selector property. This fixes $B=\binom{F+t}{t}$.
\item Form the finite weighted pending-state path. Choose the odd multiplier
of $F$ so $y$ bounds this path and every prime at most $t$ divides $y$.
Finite CRT chooses an arbitrarily large positive beta code $x$ satisfying the
state residues and endpoints. Put $N=B^x$.
\item The generalized binomial CRT supplies the nine representative classes
modulo $N^2$. Choose the eight ordinary representatives positive first.
Because $N$ is odd, the ninth class has cofinally many even representatives.
Choose $C$ even and larger than every other nominal index and representative.
The prescribed index class gives $B\mid C+F+1$.
\item Set $D=2(C+F+1)$ and $q=D^x$ only now. Then $N\mid q$ and $D>2C$.
Evaluate the true requested power outputs, then choose the main square frame
and the positive canonical main Pell pairs after all requested atoms have
been fixed. The inherited cofinal anchor
completion supplies the private comparison witnesses, orientation quotient
and cap slack with one fewer principal pair.
\item Take the normalized low remainders. The estimates of
Section~\ref{sec:radix} satisfy the retained common norm budget. Every
nonnegative slack has the inherited representation
$r^2+s^2+w(w+1)$, supplying the remaining integer witnesses.
\end{enumerate}
This constructs a zero of the changed arithmetic source from every finite
refutation. No guessed universal-machine parameter or unknown constant is
required.

\subsection{The precise logical dependency}

The repository's logical endpoints are the named theorems
\texttt{KEndpoint.c222\_iff\_zf\_refutation} and
\texttt{ConstructibleChoice.zf\_refutation\_iff\_zfc\_refutation}.
The first is a theorem about the numerical calculus itself and has no machine
argument. The original square-shell recoding transports that calculus without
changing finite refutability. The second transfers finite refutations between
ZF and ZFC. Combining these with the two arithmetic directions proves
Theorem~\ref{thm:main}.

The inherited model argument is internal to the background set theory; it
does not replace consistency by the stronger existence of an externally
well-founded transitive model. Selected logical files and the exact ordinary
proof dependency audit are preserved in the frozen parent. That collection
is not a complete compilable Lean dependency closure. This report supplies
ordinary proofs formalizable in ZFC and exact finite artifact verification,
not a new end-to-end proof-assistant build.

\section{Exact projection, coefficients and reproducibility}\label{sec:finite}

The unoptimized changed source has 216 coordinates and 156 quadratic rows.
Five certified common-Hessian projections and one affine projection give
210 coordinates and 150 rows. Their use here is entirely algebraic.

\begin{lemma}[Integral common-Hessian projection]
Let $Q_i$ be integer quadratic polynomials and let an integer direction $v$
lie in the kernel of every Hessian. Suppose $v_p=1$ and some directional
linear coefficient $\ell_r$ is $\epsilon\in\{1,-1\}$. Write
$x=y+sv$ with $y_p=0$. Then $Q_i(x)=Q_i(y)+\ell_i s$, and eliminating
$s=-\epsilon Q_r(y)$ gives the quadratic rows
$Q_i(y)-\epsilon\ell_iQ_r(y)$ for $i\ne r$. This is a bijection of the
integer zero sets, with an explicit integral reconstruction map.
\end{lemma}
The proof is the displayed polynomial identity; no rational solvability or
divisibility side condition is introduced. The affine projection uses the
actual relation
\[
4+12x_8+16x_9+4x_{19}-x_{20}+x_{33}=0
\]
at the 211-coordinate stage. It is the difference of two residuals and has
unit slope in the eliminated coordinate.

Finally 79 elementary integer operations are applied: coordinate shears or
translations with unit diagonal, and additions of integral multiples of
one residual to another. Each operation has an integer inverse. Changing
the residual basis may change the sum-of-squares polynomial, so polynomial
identity with the predecessor is not asserted. Its integer zero set is
preserved, which is the required assertion. The independent checker verifies
all forward residual identities and their composed inverse.

\begin{center}\small
\begin{tabular}{lrrr}
\toprule Stage & Coordinates & Quadratic rows & Collected monomials\\
\midrule Frozen source &219&158&\multicolumn{1}{c}{in parent}\\
Standalone Pell growth &217&157&6,943\\
Affine radix composition &216&156&6,972\\
Five Hessian projections &211&151&6,926\\
One affine projection &210&150&6,975\\
Integral charts and row operations &210&150&5,282\\
\bottomrule
\end{tabular}
\end{center}

The coordinate and row counts are literal. All final coordinates occur in
the collected polynomial. The archived JSON contains each residual as a
list of nonzero integer coefficients and sorted coordinate-index multisets.
A second JSON contains every collected coefficient of the quartic. The
principal equation JSON has SHA-256
\begin{quote}\ttfamily\small
20a2108353a3756ed7d1ddfdc71695131dcf0b6d94039988d9565ac48e6e1f0a.
\end{quote}

The supplied verifier reconstructs the source changes independently of the
producer, checks the entire main Pell interface, checks the projections,
checks all integer operations and the complete expansion, and rejects 22
deliberate corruptions of the direct construction. The printed residual
appendix is generated from the frozen final equation and checked back against
it. Tests also run with Python's optimization flag, so substantive checks do
not depend on \texttt{assert} statements.

The commands \texttt{python3 -B verify.py} and
\texttt{python3 -B -O verify.py} use only the Python standard library and
need no network. To reconstruct the coefficients in a fresh directory, run
\begin{quote}\small\ttfamily
python3 -B code/rebuild.py\par
\hspace{1em}--output-dir /tmp/conzfc-210-rebuilt
\end{quote}
as one command.
Finite checks authenticate polynomial coefficients and transformations; the
unbounded Pell, CRT and logical arguments are the mathematical proofs above
and the explicitly preserved dependencies.

\section{An independent compact companion}\label{sec:compact}

\label{sec:constant-dilation}
The direct quartic and the compact construction in this section have different
coefficient architectures. Here the coefficients are the complete retained
BFS943/radix67 tuple $(\beta,K_0,M_C,M_F)$, whose largest member has
$592884$ bits. The new equation has eleven strictly positive integer
unknowns, exact degree $174$, and $104$ live binary arithmetic operations.
It retains the native/compiler/refutation calibration. Its full ordinary
proof and all coefficient bytes are included in the companion directory.
No equality with the direct quartic is asserted.

The inspected source-115 comparator has eleven positive unknowns, degree
$211$, and $107$ operations. Its four coefficient bit lengths agree but its
whole-tuple $K_0$ differs. Thus the comparison concerns domain and size,
not an identity between the specialized polynomials. The source-172 equation
with ten positive unknowns, degree $214$, and $87$ operations uses the same
complete tuple as the present construction and remains a distinct tradeoff.

\subsection{The complete polynomial}
All symbols in this section are local. The independent inputs are
\[
 J,F,\alpha,i,T,w,Y,n,\eta,\zeta,\kappa\in\mathbb Z_{>0}.
\]
The circuit JSON retains the historical input name \texttt{s} for $Y$;
this is a renaming of one input, not an additional variable.
Every other symbol below is a polynomial abbreviation. Set $\lambda=210$ and
\begin{align*}
 h&=\beta J,& q&=h+1,\\
 C&=q-F-\alpha,&X&=qw,\\
 k&=\eta+\zeta,&c&=kY+\eta,\\
 A&=\lambda Y(X+1)+1,&\Delta&=A^2-1,\\
 H&=4A-5,&v_0&=XY^2,\\
 v&=\lambda^2v_0,&\ell&=vk,\\
 E&=1+\ell(\ell+k),&B&=(K_0+w)C-F.
\end{align*}
\[
 R=[h(q-F)+\alpha](q^2-1)+(M_C+qM_F)J.
\]
\begin{align*}
 U&=ic^2,&a&=\Delta U^2E,\\
 F_0&=1+a,&G&=EF_0,\\
 f&=1+2a,&Q&=4\Delta aF_0,\\
 b&=cT-R,&V&=fb-c,\\
 L&=bV,&K&=X(3c-2X)+2.
\end{align*}
Define the fixed integer polynomial
\begin{align*}
 \mathcal T(k,x)={}&k-7350xk(k^2-4)\\
 &+5402250x^2k(k^2-4)(27k^2+8)\\
 &-5672362500x^3k(k^2-4)(675k^4-324k^2+160).
\end{align*}
The final registers are
\begin{align*}
 e_*&=\mathcal T(k,v_0)-R-1,&D&=q^7,&M&=hD,\\
 Z&=M(n+\kappa-LB)-DB-h(e_*-q^4Y),\\
 P_0&=Q(Z^2-M^2GV^2)+(M^2G-Z^2),&N&=K-\kappa H.
\end{align*}
The equation is
\[
 \boxed{\quad P=(P_0+1)(3N+1)-1=0.\quad}
\]
There are no extra quantified registers, radicals, variable exponents or
division gates. Since $P_0,N$ are integers, $P=0$ forces both displayed
factors to be units. The second factor is $1$ modulo $3$, and therefore
cannot be $-1$. Consequently
\[
 P=0\quad\Longleftrightarrow\quad P_0=N=0.
\]
The independently integral projector quotient $\kappa$ is essential.

\subsection{An inverse recurrence with integral coefficients}
Let $\psi_b(0)=0$, $\psi_b(1)=1$ and
$\psi_b(m+2)=2b\psi_b(m+1)-\psi_b(m)$. Direct recurrence induction gives
\[
 2\psi_{1+2v}(m)
 =2\sum_{j=0}^{m-1}\binom{m+j}{2j+1}(4v)^j.
\]
Put
\begin{align*}
 C_1(k)&=-\frac{k(k^2-4)}6,\\
 C_2(k)&=\frac{k(k^2-4)(27k^2+8)}{360},\\
 C_3(k)&=-\frac{k(k^2-4)(675k^4-324k^2+160)}{15120}.
\end{align*}
These are the first three corrections in the formal inverse
$k+C_1(k)v+C_2(k)v^2+C_3(k)v^3$. With $\lambda=210$, their denominators
clear before any congruence argument:
\[
 \lambda^2/6=7350,\qquad
 \lambda^4/360=5402250,\qquad
 \lambda^6/15120=5672362500.
\]
For every integer $m\ge1$ the following is divisibility in
$\mathbb Z[x]$:
\begin{equation}\label{eq:inverse-third}
 \mathcal T\bigl(2\psi_{1+2\lambda^2x}(m),x\bigr)-2m
 \in x^4\mathbb Z[x].
\end{equation}
Indeed, the first four recurrence coefficients are
\[
 2m+8\lambda^2\binom{m+1}{3}x
 +32\lambda^4\binom{m+2}{5}x^2
 +128\lambda^6\binom{m+3}{7}x^3.
\]
Substitution makes the coefficients of $x,x^2,x^3$ vanish as exact
polynomials in $\mathbb Q[m]$. Both full polynomials have integer
coefficients for each integer $m$, proving \eqref{eq:inverse-third}.
The supplied finite certificate checks these polynomial identities; it
does not replace the all-index recurrence proof by sampled values.

The transport argument needs a sign bound valid before index recovery.
For $k\ge2$ and $x\ge1$,
\begin{equation}\label{eq:inverse-third-sign}
 \mathcal T(k,x)\le k.
\end{equation}
To verify it, set $t=k^2\ge4$ and $v=\lambda^2x\ge1$. The expression
subtracted from $k$ is $k(k^2-4)$ times
\[
 \frac v6-\frac{v^2(27t+8)}{360}
 +\frac{v^3(675t^2-324t+160)}{15120}.
\]
The last two terms have nonnegative sum because
\[
 (675t^2-324t+160)-42(27t+8)
 =675(t-4)^2+3942(t-4)+4792>0.
\]

\subsection{Soundness in the required order}
The complete proof is retained in
\texttt{compact11/PROOF\_constant\_dilation.md}; the key dependencies are given here.
At $P_0=0$, write $z=Z/M$. Prime valuations in
\[
 (Q-1)z^2=G(QV^2-1)
\]
give positive integral roots $E=\tau^2$, $F_0=f_0^2$ and
$G=g^2$ with $g=\tau f_0$. This uses $Q=4\Delta^2U^2G$ and
$\gcd(E,F_0)=1$. Adjacent-coefficient Pell descent then gives
$y=|z|/g\in\mathbb Z_{>0}$, so $z=\pm gy$ is integral.
The explicit strong pair has
\[
 I=2Ug,\qquad f^2-\Delta I^2=1,\qquad Q=(\Delta I)^2,\qquad c^2\mid I.
\]
The retained signed all-index folding lemma yields
$c=\psi_A(p)$ for an odd $p\ge3$. The first norm yields
$k=2\psi_{1+2v}(m)$. Its index is at least two: $m=1$ would give
$k=2$ and $c=2Y+1$ strictly between the first two positive main Pell
ordinates. Since $1+2v>A$ and $c>k$, one also has $m<p$.

The free chart no longer has the parent's automatic bound $A>q^2$.
A separate argument supplies the needed replacement $c>2q^4$.
If $p\ge5$, then $c\ge16A^4-12A^2+1>4A^4>4q^4$.
If $p=3$, then $m=2$ and
\[
 k=8\lambda^2XY^2+4,\qquad c=4A^2-1,\qquad
 \frac ck>\frac{(X+1)^2}{2X+1}>\frac{X+1}{2}.
\]
Together with $c/k<Y+1$, integrality gives $2Y\ge X$, whence
$c>\lambda^2X^4-1>2q^4$.

Before the sign of $R$ is known, the uniform auxiliary estimate gives
\[
 L=bV>0,\qquad gy<3L/8,\qquad |b|/q^7<L/2.
\]
For completeness, $b=0$ would imply
$0<y^2-c^2=(c^2-1)/(Q-1)<1$. The auxiliary Pell index is odd and at
least three, so $|V|\ge4Q-3$. With $c_3=4\Delta+3\le c$ and $f>2c$,
\[
 |b|\ge4\Delta f-\frac{c+c_3}{f}>4\Delta f-1>4g,
\]
while $y<3|V|/2$. These prove the estimates.

Use the proof-only quantities
\[
 \sigma=Y/q^3,\quad \rho=e_*/q^7,\quad W=1+hL,\quad t=B/h.
\]
The carrier identity is
\[
 z=n+\kappa+\sigma-Wt-\rho.
\]
By \eqref{eq:inverse-third-sign}, $e_*\le k-R-1<b$.
If $B\le-1$, the identity gives
$n<L/2-L+|b|/q^7<0$. Thus $B\ge0$ and $C>0$.
The retained mask estimate now gives $3q<R<q^4-2$.
Hence $c>2R$, and the normalized strong-rank and signed auxiliary-index
lemmas apply to $V\equiv-c\pmod f$ and $V\equiv-R\pmod c$.
They give $p=R$.

Integral $z$ implies $q^7\mid h(e_*-q^4Y)$. Since $\gcd(q,h)=1$,
$q^4\mid e_*$. Equation \eqref{eq:inverse-third}, together with $q\mid v_0$,
gives $q^4\mid e_*-(2m-R-1)$. But $1\le m<R<q^4$, and therefore
\[
 2m=R+1.
\]
No coprimality between $q$ and $H$ has been used.

Write $R=2r+1$, so $r\ge48$. The inherited ratio theorem is uniform for
every $\lambda\ge1$, $Y\ge1$:
\[
 \frac{(X+1)^{2r}}{X^r}
 <\frac{\psi_A(2r+1)}{\psi_{1+2v}(r+1)}
 <\frac{(X+1)^{2r}}{X^r}
      \left(1+\frac1{Y(X+1)}\right)^{2r}.
\]
It implies $Y>X^r/3$ and $H>X^{r+1}$.
The independent equation $N=0$ gives $H\mid K$; the universal Pell
projector then gives
\[
 H\mid (X-2^R)(2^RX+1).
\]
The entire product has absolute value below
$4\cdot4^rX^2+8\cdot16^rX<X^{r+1}<H$, so it vanishes and $X=2^R$.

Since $q\mid X$, write $q=2^a$ with $a\ge5$. The inequality
$R>3q\ge7a$ gives $q^7\mid X$, hence $q^7\mid v_0$.
Equation \eqref{eq:inverse-third} now gives $q^7\mid e_*$.
Thus $\rho$ is integral, and the carrier gives
$Y/q^3-WB/h\in\mathbb Z$. Coprimality and $W\equiv1\pmod h$ yield
\[
 q^3\mid Y,\qquad h\mid B,\qquad B/h\ge0.
\]
The retained fixed-minus parity, half-binomial rounding and fresh native
reconstruction apply. At the entire supplied fixed tuple they produce a
finite syntactic ZFC refutation.

\subsection{Positive completeness and exact size}
Conversely, a refutation supplies the retained canonical outer data,
including $X=2^R$, $R\equiv3\pmod4$, $q^3\mid Y$ and $B/h\ge0$.
Form the new bases with $\lambda=210$ and set
\[
 c=\psi_A(R),\quad k=2\psi_{1+2v}((R+1)/2),\quad
 \tau=\chi_{1+2v}((R+1)/2).
\]
The uniform ratio theorem gives positive integers
$\eta=c-kY$ and $\zeta=k-\eta$. Projection gives positive integral
$\kappa=K/H$, while \eqref{eq:inverse-third} gives integral $\rho$.

A positive period of the determinant-one Pell matrix modulo $c^2\tau$
provides arbitrarily large pairs $(f_0,I_0)$ with $c^2\tau\mid I_0$.
Take $i=I_0/(c^2\tau)$, so $g=\tau f_0$, and choose the period multiple
large enough that $g>\kappa+\sigma+|\rho|$.
The doubled pair gives $f=2f_0^2-1$, $I=2f_0I_0$ and $S=\Delta I$.
Put
\[
 V=\chi_S(R)/S,\qquad y=\psi_S(R),\qquad
 T=\frac{V+c+Rf}{cf}.
\]
The odd quotient identities at $R\equiv3\pmod4$ make $T$ a positive
integer and give $V=f(cT-R)-c$. Finally
\[
 n=gy+(1+hL)(B/h)+\rho-\kappa-\sigma>0.
\]
These witnesses satisfy both $P_0=0$ and $N=0$.
The cofinal choice of the strong period is essential because the new
inverse correction may be negative and is larger than the old correction.

The complete highest homogeneous part, with $q_0=\beta J$ and
$k_0=\eta+\zeta$, is
\[
 [P]_{174}=
 -48\lambda^{36}q_0^{47}w^{25}Y^{56}i^{10}k_0^{32}(q_0-F)^2
 \bigl(3k_0Y-2q_0w-4\lambda\kappa Y\bigr).
\]
In particular,
\[
 [J^{49}w^{25}Y^{57}i^{10}\eta^{32}\kappa]P
 =192\lambda^{37}\beta^{49}>0.
\]
This proves exact degree $174$ at the actual fixed coefficients.
The independent checker verifies all $104$ live gates by exact polynomial
cuts, the inverse identity in $\mathbb Q[m,x]/(x^4)$, the complete
$300$-monomial leading form, all four reconstructed coefficient payloads,
and rejection of $208$ operator mutations. These finite checks support
the emitted artifacts; the preceding unbounded argument supplies the
zero-existence equivalence.

Deleting $\kappa$ and restoring $K/H$ inside the shared carrier would
suggest a ten-variable degree-$176$ formula. That shortcut is rejected:
fixed $\lambda$ and free $Y$ do not guarantee $\gcd(q,H)=1$, so its first-index
recovery is not proved.



\section{Further research questions}\label{sec:further}

\paragraph{Compress the remaining principal-power interface.}
Thirteen principal main-Pell pairs remain. The new result shows why a special
purpose growth certificate can be cheaper than an exact-power certificate.
An analogous specialization of one remaining request is a substantive target.
Deleting a pair merely because its nominal index is related to another index
would require a proved replacement for all orientation and digit obligations.

\paragraph{Reduce the growth multiplier without increasing degree.}
The coefficient $2^{19}$ arises from the depth-19 syntax majorant and a
2-adic index constraint. A shorter equivalent seed term, a tighter
source-specific polynomial majorant, or a different congruence forcing a
larger index per coefficient could lower height. The current depth argument
is sufficient, not a lower bound on the best possible encoding.

\paragraph{Eliminate a coordinate beyond constant Hessian directions.}
The inherited projections exploit common Hessian kernels and constant
residual-row combinations. Nonconstant polynomial combinations, a different
Pell chart, or a structural change in the local calculus could expose a new
integral projection. Any proposed elimination must retain integer
reconstruction and report the resulting support and coefficient growth.

\paragraph{Remove the compact companion's separate projector quotient.}
The tempting ten-variable shortcut is blocked by a coprimality requirement:
the fixed dilation no longer automatically gives $\gcd(q,H)=1$.
A new carrier that proves the necessary denominator separation could improve
the compact frontier. Merely substituting a rational quotient into the
eleven-variable construction is not justified.

\paragraph{Build a complete formal proof chain.}
The most valuable verification extension is to formalize the new standalone
growth lemma, finite binary-residue lifting, prime-support beta CRT, and
changed-source decoding, then compile the full dependency closure through
the numerical-calculus/ZF/ZFC endpoints. More bounded examples alone cannot
close this formalization boundary.

\paragraph{Specify an optimization criterion.}
The present 210-variable equation and the 214-variable small-height companion
are useful different outcomes. Future searches should state whether the
priority is arity at degree four, coefficient height, collected support,
binary circuit length, or a weighted combination. No proof here excludes
a dramatically smaller equation in any of these senses.

\begingroup\raggedright
\begin{thebibliography}{9}
\bibitem{repo} V. Reshetnikov, \emph{ProveIt}, inspected snapshot
\texttt{ef0c534c1ae1e312c61a2a9a5a3f18f5516fd42b},
\url{https://github.com/VladimirReshetnikov/ProveIt}.
Relevant directories: \texttt{Computability/HilbertTenthProblem},
\texttt{Computability/DiophantineConsistency/con-zfc-diophantine-forms},
and \texttt{docs/incoming}.
\bibitem{anchor} \emph{A reused CRT index: an explicit Diophantine consistency
equation}, incoming research package, 25 October 2026,
\texttt{ConZFC\_CRT\_Index\_Anchor\_215\_2026-10-25.zip};
archive SHA-256 \texttt{760428ab91a1f39019ebb452d8b347235b5bb1a0935346e23de78b4542e68677}.
Its inherited source snapshot is \texttt{953fb4d82d4229f4efc5f9db4cc2bf38ae329bb4}.
\bibitem{matiyasevich} Yu. V. Matiyasevich,
\emph{Diophantine representation of enumerable predicates},
Izv. Akad. Nauk SSSR Ser. Mat. \textbf{35}(1), 3--30 (1971);
English translation, Mathematics of the USSR-Izvestiya \textbf{5}, 1--28.
\url{https://www.mathnet.ru/eng/im1910}.
\bibitem{compactsource} ProveIt entries 115 and 172, compact equation
comparison and retained source; exact file paths, byte hashes, coefficient
tuple and inherited proofs are included in \texttt{compact11/}.
\end{thebibliography}
\endgroup

\clearpage
\appendix
\section{The complete literal principal equation}\label{app:literal}

Every coordinate $x_0,\ldots,x_{209}$ in this appendix is an unrestricted
integer. The following 150 displayed residuals define the equation without
any omitted coefficient, schematic index, exponentiation request, or external
machine parameter:
\[
\boxed{\ P_{210}=R_0^2+R_1^2+\cdots+R_{149}^2=0.\ }
\]
The sum-of-squares form is itself an explicit polynomial expression. Its
5,282 collected monomials are additionally supplied in
\texttt{equations/principal210/expanded\_polynomial.json}.
The final three rows display the new growth block directly: $t=x_0$,
$F=x_2$, $X=x_{207}$, $Z=x_{208}$ and $u=x_{209}$.

\begingroup\small
% Literal residuals; checked independently against equation.json.
{\small
\setlength{\jot}{2pt}
\allowdisplaybreaks[0]
% BEGIN RESIDUAL 0
\begin{align*}
R_{0} ={}& x_{0} - x_{0} x_{4} - x_{4} + x_{7}
\end{align*}
% END RESIDUAL 0
% BEGIN RESIDUAL 1
\begin{align*}
R_{1} ={}& -x_{2} - x_{4} x_{5} + x_{6} x_{7}
\end{align*}
% END RESIDUAL 1
% BEGIN RESIDUAL 2
\begin{align*}
R_{2} ={}& -1 + x_{1} - x_{6} x_{7}
\end{align*}
% END RESIDUAL 2
% BEGIN RESIDUAL 3
\begin{align*}
R_{3} ={}& -2x_{2} x_{206} - x_{3} x_{18} + x_{4} + 2x_{16}
\end{align*}
% END RESIDUAL 3
% BEGIN RESIDUAL 4
\begin{align*}
R_{4} ={}& -4 - 3x_{8} - x_{8}^{2} - 2x_{8} x_{9} - 4x_{9} - x_{9}^{2} \\
&{}  + x_{19}
\end{align*}
% END RESIDUAL 4
% BEGIN RESIDUAL 5
\begin{align*}
R_{5} ={}& x_{8} - x_{8}^{2} - 8x_{8} x_{19} - 4x_{8} x_{20} - 16x_{19}^{2} - 16x_{19} x_{20} \\
&{}  - 4x_{20}^{2} + x_{21}
\end{align*}
% END RESIDUAL 5
% BEGIN RESIDUAL 6
\begin{align*}
R_{6} ={}& 1 + 2x_{21} - 4x_{21}^{2} + x_{22}
\end{align*}
% END RESIDUAL 6
% BEGIN RESIDUAL 7
\begin{align*}
R_{7} ={}& -2 - x_{9}^{2} - 2x_{9} x_{10} - x_{10}^{2} + x_{23}
\end{align*}
% END RESIDUAL 7
% BEGIN RESIDUAL 8
\begin{align*}
R_{8} ={}& 1 - x_{8}^{2} - 2x_{8} x_{10} - x_{10}^{2} + x_{24}
\end{align*}
% END RESIDUAL 8
% BEGIN RESIDUAL 9
\begin{align*}
R_{9} ={}& -29 - 61x_{8} - 35x_{8}^{2} - 96x_{8} x_{10} - 24x_{8} x_{19} - 4x_{8} x_{23} \\
&{}  - 24x_{8} x_{24} + 16x_{9} - 2x_{9}^{2} - 8x_{9} x_{23} - 96x_{10} - 64x_{10}^{2} \\
&{}  - 32x_{10} x_{19} - 32x_{10} x_{24} - 28x_{19} - 4x_{19}^{2} - 8x_{19} x_{24} + 8x_{23} \\
&{}  - 4x_{23}^{2} - 24x_{24} - 4x_{24}^{2} + x_{25}
\end{align*}
% END RESIDUAL 9
% BEGIN RESIDUAL 10
\begin{align*}
R_{10} ={}& -4 + 10x_{8} - 2x_{8}^{2} - 8x_{8} x_{9} - 8x_{8} x_{23} - 4x_{25}^{2} \\
&{}  + x_{26} - 2x_{43}
\end{align*}
% END RESIDUAL 10
% BEGIN RESIDUAL 11
\begin{align*}
R_{11} ={}& -4 - 16x_{10}^{2} - 16x_{10} x_{23} - 16x_{10} x_{24} - 4x_{23}^{2} - 8x_{23} x_{24} \\
&{}  - 4x_{24}^{2} + x_{27}
\end{align*}
% END RESIDUAL 11
% BEGIN RESIDUAL 12
\begin{align*}
R_{12} ={}& -7 + 24x_{10} - 16x_{10}^{2} - 16x_{10} x_{19} - 32x_{10} x_{24} - 16x_{10} x_{27} \\
&{}  + 10x_{19} - 4x_{19}^{2} - 16x_{19} x_{24} - 8x_{19} x_{27} + 24x_{24} - 16x_{24}^{2} \\
&{}  - 16x_{24} x_{27} + 12x_{27} - 4x_{27}^{2} + x_{28}
\end{align*}
% END RESIDUAL 12
% BEGIN RESIDUAL 13
\begin{align*}
R_{13} ={}& -114 - 127x_{8} - 37x_{8}^{2} - 24x_{8} x_{9} - 112x_{8} x_{10} - 4x_{8} x_{23} \\
&{}  - 24x_{8} x_{24} - 67x_{9} - 37x_{9}^{2} - 112x_{9} x_{10} - 24x_{9} x_{23} - 4x_{9} x_{24} \\
&{}  - 224x_{10} - 128x_{10}^{2} - 32x_{10} x_{23} - 32x_{10} x_{24} - 16x_{23} - 4x_{23}^{2} \\
&{}  - 40x_{24} - 4x_{24}^{2} + x_{29}
\end{align*}
% END RESIDUAL 13
% BEGIN RESIDUAL 14
\begin{align*}
R_{14} ={}& -30 - 14x_{8} - 2x_{8}^{2} - 24x_{8} x_{9} - 32x_{8} x_{10} - 8x_{8} x_{23} \\
&{}  - 96x_{9} - 72x_{9}^{2} - 192x_{9} x_{10} - 48x_{9} x_{23} - 128x_{10} - 128x_{10}^{2} \\
&{}  - 64x_{10} x_{23} - 32x_{23} - 8x_{23}^{2} - 4x_{29}^{2} + x_{30}
\end{align*}
% END RESIDUAL 14
% BEGIN RESIDUAL 15
\begin{align*}
R_{15} ={}& -74 - 42x_{8} + 8x_{8} x_{10} + 16x_{8} x_{23} - 84x_{9} + 16x_{9} x_{10} \\
&{}  + 32x_{9} x_{23} + 29x_{10} - 3x_{10}^{2} - 4x_{10} x_{19} - 8x_{10} x_{23} + 21x_{19} \\
&{}  - 8x_{19} x_{23} + 40x_{23} - 4x_{23}^{2} + x_{31} - x_{59}
\end{align*}
% END RESIDUAL 15
% BEGIN RESIDUAL 16
\begin{align*}
R_{16} ={}& -131 - 184x_{8} - 64x_{8}^{2} - 160x_{8} x_{10} - 64x_{8} x_{23} - 32x_{8} x_{24} \\
&{}  - 32x_{8} x_{31} - 226x_{10} - 100x_{10}^{2} - 80x_{10} x_{23} - 40x_{10} x_{24} - 40x_{10} x_{31} \\
&{}  - 88x_{23} - 16x_{23}^{2} - 16x_{23} x_{24} - 16x_{23} x_{31} - 46x_{24} - 4x_{24}^{2} \\
&{}  - 8x_{24} x_{31} - 44x_{31} - 4x_{31}^{2} + x_{32}
\end{align*}
% END RESIDUAL 16
% BEGIN RESIDUAL 17
\begin{align*}
R_{17} ={}& -4x_{19}^{2} + x_{20}
\end{align*}
% END RESIDUAL 17
% BEGIN RESIDUAL 18
\begin{align*}
R_{18} ={}& x_{8} - x_{8}^{2} + 8x_{8} x_{19} - 4x_{8} x_{20} - 16x_{19}^{2} + 16x_{19} x_{20} \\
&{}  - 4x_{20}^{2} + x_{33}
\end{align*}
% END RESIDUAL 18
% BEGIN RESIDUAL 19
\begin{align*}
R_{19} ={}& 1 + 2x_{33} - 4x_{33}^{2} + x_{34}
\end{align*}
% END RESIDUAL 19
% BEGIN RESIDUAL 20
\begin{align*}
R_{20} ={}& -x_{9}^{2} - 4x_{9} x_{34} - 4x_{34}^{2} + x_{35}
\end{align*}
% END RESIDUAL 20
% BEGIN RESIDUAL 21
\begin{align*}
R_{21} ={}& -44 - 52x_{8} - 16x_{8}^{2} - 64x_{8} x_{9} + 32x_{8} x_{19} - 104x_{9} \\
&{}  - 64x_{9}^{2} + 64x_{9} x_{19} + 50x_{19} - 16x_{19}^{2} + x_{36}
\end{align*}
% END RESIDUAL 21
% BEGIN RESIDUAL 22
\begin{align*}
R_{22} ={}& 8x_{9} - 62x_{9}^{2} - 92x_{9} x_{10} - 32x_{9} x_{23} - 32x_{9} x_{36} + 6x_{10} \\
&{}  - 34x_{10}^{2} - 24x_{10} x_{23} - 24x_{10} x_{36} - 4x_{23}^{2} - 8x_{23} x_{36} + 4x_{36} \\
&{}  - 4x_{36}^{2} + x_{37}
\end{align*}
% END RESIDUAL 22
% BEGIN RESIDUAL 23
\begin{align*}
R_{23} ={}& -90 - 152x_{8} - 64x_{8}^{2} - 96x_{8} x_{10} - 32x_{8} x_{24} - 32x_{8} x_{37} \\
&{}  - 114x_{10} - 36x_{10}^{2} - 24x_{10} x_{24} - 24x_{10} x_{37} - 38x_{24} - 4x_{24}^{2} \\
&{}  - 8x_{24} x_{37} - 36x_{37} - 4x_{37}^{2} + x_{38}
\end{align*}
% END RESIDUAL 23
% BEGIN RESIDUAL 24
\begin{align*}
R_{24} ={}& -x_{10}^{2} - 2x_{10} x_{11} - x_{11}^{2} + x_{39}
\end{align*}
% END RESIDUAL 24
% BEGIN RESIDUAL 25
\begin{align*}
R_{25} ={}& -7 + 8x_{8} - 3x_{8}^{2} - 8x_{8} x_{9} - 8x_{8} x_{23} - 8x_{8} x_{24} \\
&{}  - 4x_{8} x_{39} - 8x_{9} x_{24} - 2x_{10}^{2} - 8x_{10} x_{39} - 8x_{23} x_{24} + 6x_{24} \\
&{}  - 4x_{24}^{2} - 4x_{39}^{2} + x_{40} - x_{43}
\end{align*}
% END RESIDUAL 25
% BEGIN RESIDUAL 26
\begin{align*}
R_{26} ={}& -16x_{8}^{2} - 36x_{8} x_{10} - 32x_{8} x_{24} - 40x_{8} x_{39} - 16x_{8} x_{40} - 4x_{10} \\
&{}  - 22x_{10}^{2} - 32x_{10} x_{24} - 48x_{10} x_{39} - 16x_{10} x_{40} - 2x_{24} - 16x_{24}^{2} \\
&{}  - 32x_{24} x_{39} - 16x_{24} x_{40} - 4x_{39} - 24x_{39}^{2} - 16x_{39} x_{40} - 4x_{40}^{2} \\
&{}  + x_{41}
\end{align*}
% END RESIDUAL 26
% BEGIN RESIDUAL 27
\begin{align*}
R_{27} ={}& x_{10} - 25x_{10}^{2} - 40x_{10} x_{24} - 20x_{10} x_{27} - 16x_{24}^{2} - 16x_{24} x_{27} \\
&{}  - 4x_{27}^{2} + x_{42}
\end{align*}
% END RESIDUAL 27
% BEGIN RESIDUAL 28
\begin{align*}
R_{28} ={}& -1 + 8x_{9} - 4x_{9}^{2} - 8x_{9} x_{23} + 8x_{23} - 4x_{23}^{2} \\
&{}  + x_{43}
\end{align*}
% END RESIDUAL 28
% BEGIN RESIDUAL 29
\begin{align*}
R_{29} ={}& -1 - 16x_{9}^{2} - 8x_{9} x_{10} - 32x_{9} x_{23} - 16x_{9} x_{43} - x_{10}^{2} \\
&{}  - 8x_{10} x_{23} - 4x_{10} x_{43} - 16x_{23}^{2} - 16x_{23} x_{43} - 4x_{43}^{2} + x_{44}
\end{align*}
% END RESIDUAL 29
% BEGIN RESIDUAL 30
\begin{align*}
R_{30} ={}& 3 + x_{9} - 2x_{9} x_{10} - 4x_{9} x_{44} - 3x_{10}^{2} - 8x_{10} x_{44} \\
&{}  - x_{23} - 4x_{44}^{2} + x_{45}
\end{align*}
% END RESIDUAL 30
% BEGIN RESIDUAL 31
\begin{align*}
R_{31} ={}& -1 - 4x_{45} - 4x_{45}^{2} + x_{46}
\end{align*}
% END RESIDUAL 31
% BEGIN RESIDUAL 32
\begin{align*}
R_{32} ={}& -1 - x_{10}^{2} - 8x_{10} x_{45} - 4x_{10} x_{46} - 16x_{45}^{2} - 16x_{45} x_{46} \\
&{}  - 4x_{46}^{2} + x_{47}
\end{align*}
% END RESIDUAL 32
% BEGIN RESIDUAL 33
\begin{align*}
R_{33} ={}& 2 + x_{9} - x_{9}^{2} - 4x_{9} x_{10} - 4x_{9} x_{47} - 4x_{10}^{2} \\
&{}  - 8x_{10} x_{47} - 4x_{47}^{2} + x_{48}
\end{align*}
% END RESIDUAL 33
% BEGIN RESIDUAL 34
\begin{align*}
R_{34} ={}& 1 - 4x_{42}^{2} - 8x_{42} x_{48} - 4x_{48}^{2} + x_{49}
\end{align*}
% END RESIDUAL 34
% BEGIN RESIDUAL 35
\begin{align*}
R_{35} ={}& 1 + x_{8} - x_{8}^{2} - 8x_{8} x_{42} - 4x_{8} x_{49} - 16x_{42}^{2} \\
&{}  - 16x_{42} x_{49} - 4x_{49}^{2} + x_{50}
\end{align*}
% END RESIDUAL 35
% BEGIN RESIDUAL 36
\begin{align*}
R_{36} ={}& -1 - 4x_{50} - 4x_{50}^{2} + x_{51}
\end{align*}
% END RESIDUAL 36
% BEGIN RESIDUAL 37
\begin{align*}
R_{37} ={}& -1 - x_{10}^{2} - 8x_{10} x_{50} - 4x_{10} x_{51} - 16x_{50}^{2} - 16x_{50} x_{51} \\
&{}  - 4x_{51}^{2} + x_{52}
\end{align*}
% END RESIDUAL 37
% BEGIN RESIDUAL 38
\begin{align*}
R_{38} ={}& 2 + x_{8} - x_{8}^{2} - 4x_{8} x_{10} - 4x_{8} x_{52} - 4x_{10}^{2} \\
&{}  - 8x_{10} x_{52} - 4x_{52}^{2} + x_{53}
\end{align*}
% END RESIDUAL 38
% BEGIN RESIDUAL 39
\begin{align*}
R_{39} ={}& -10 - 12x_{53} - 4x_{53}^{2} + x_{54}
\end{align*}
% END RESIDUAL 39
% BEGIN RESIDUAL 40
\begin{align*}
R_{40} ={}& -19 + 4x_{9} - 4x_{9}^{2} - 8x_{9} x_{23} - 16x_{9} x_{53} - 8x_{9} x_{54} \\
&{}  + 4x_{23} - 4x_{23}^{2} - 16x_{23} x_{53} - 8x_{23} x_{54} - 16x_{53} - 24x_{53}^{2} \\
&{}  - 16x_{53} x_{54} + 6x_{54} - 4x_{54}^{2} + x_{55}
\end{align*}
% END RESIDUAL 40
% BEGIN RESIDUAL 41
\begin{align*}
R_{41} ={}& -10 + x_{9} - 25x_{9}^{2} - 40x_{9} x_{23} - 20x_{9} x_{55} - 16x_{23}^{2} \\
&{}  - 16x_{23} x_{55} - 12x_{48} - 4x_{48}^{2} - 4x_{55}^{2} + x_{56}
\end{align*}
% END RESIDUAL 41
% BEGIN RESIDUAL 42
\begin{align*}
R_{42} ={}& 21 + 24x_{48} - 8x_{48}^{2} - 16x_{48} x_{56} - 2x_{56} - 4x_{56}^{2} \\
&{}  + x_{57}
\end{align*}
% END RESIDUAL 42
% BEGIN RESIDUAL 43
\begin{align*}
R_{43} ={}& -1 + 4x_{8} + 8x_{9} + x_{10} - x_{10}^{2} - 4x_{10} x_{57} \\
&{}  - 2x_{19} - 4x_{57}^{2} + x_{58} + 12x_{59} - 4x_{59}^{2}
\end{align*}
% END RESIDUAL 43
% BEGIN RESIDUAL 44
\begin{align*}
R_{44} ={}& -22 - 38x_{8} - 16x_{8}^{2} - 64x_{8} x_{9} + 8x_{8} x_{10} + 16x_{8} x_{19} \\
&{}  - 76x_{9} - 64x_{9}^{2} + 16x_{9} x_{10} + 32x_{9} x_{19} + 11x_{10} - x_{10}^{2} \\
&{}  - 4x_{10} x_{19} + 19x_{19} - 4x_{19}^{2} + x_{59}
\end{align*}
% END RESIDUAL 44
% BEGIN RESIDUAL 45
\begin{align*}
R_{45} ={}& -4 + 8x_{8} + 16x_{9} - 4x_{19} - 4x_{58}^{2} + 24x_{59} \\
&{}  - 8x_{59}^{2} + x_{60}
\end{align*}
% END RESIDUAL 45
% BEGIN RESIDUAL 46
\begin{align*}
R_{46} ={}& -1 - 4x_{60} - 4x_{60}^{2} + x_{61}
\end{align*}
% END RESIDUAL 46
% BEGIN RESIDUAL 47
\begin{align*}
R_{47} ={}& -x_{2} x_{13} + x_{62}
\end{align*}
% END RESIDUAL 47
% BEGIN RESIDUAL 48
\begin{align*}
R_{48} ={}& x_{12} - x_{13} - 2x_{22} x_{62} - x_{63}
\end{align*}
% END RESIDUAL 48
% BEGIN RESIDUAL 49
\begin{align*}
R_{49} ={}& x_{12} - x_{13} - 2x_{26} x_{62} - x_{64}
\end{align*}
% END RESIDUAL 49
% BEGIN RESIDUAL 50
\begin{align*}
R_{50} ={}& x_{12} - x_{13} - 2x_{28} x_{62} - 2x_{62} - x_{65}
\end{align*}
% END RESIDUAL 50
% BEGIN RESIDUAL 51
\begin{align*}
R_{51} ={}& x_{12} - x_{13} - 2x_{30} x_{62} - x_{66}
\end{align*}
% END RESIDUAL 51
% BEGIN RESIDUAL 52
\begin{align*}
R_{52} ={}& x_{12} - x_{13} - 2x_{32} x_{62} + 2x_{62} - x_{67}
\end{align*}
% END RESIDUAL 52
% BEGIN RESIDUAL 53
\begin{align*}
R_{53} ={}& x_{12} - x_{13} - 4x_{34} x_{62} - 2x_{35} x_{62} - x_{68}
\end{align*}
% END RESIDUAL 53
% BEGIN RESIDUAL 54
\begin{align*}
R_{54} ={}& x_{12} - x_{13} - 2x_{38} x_{62} - x_{69}
\end{align*}
% END RESIDUAL 54
% BEGIN RESIDUAL 55
\begin{align*}
R_{55} ={}& 2x_{28} x_{62} - 2x_{41} x_{62} - 2x_{62} + x_{65} - x_{70}
\end{align*}
% END RESIDUAL 55
% BEGIN RESIDUAL 56
\begin{align*}
R_{56} ={}& -x_{8} x_{9} + x_{8} x_{10} + x_{9} x_{10} - x_{10}^{2} + x_{71}
\end{align*}
% END RESIDUAL 56
% BEGIN RESIDUAL 57
\begin{align*}
R_{57} ={}& 1 + 2x_{11} + x_{11}^{2} - x_{71}^{2} + x_{72}
\end{align*}
% END RESIDUAL 57
% BEGIN RESIDUAL 58
\begin{align*}
R_{58} ={}& -x_{12}^{2} + 2x_{12} x_{13} + 2x_{12} x_{74} - x_{13}^{2} - 2x_{13} x_{74} - x_{72}^{2} \\
&{}  + x_{73} - x_{74}^{2}
\end{align*}
% END RESIDUAL 58
% BEGIN RESIDUAL 59
\begin{align*}
R_{59} ={}& -4x_{60} x_{62} - 2x_{61} x_{62} + x_{74}
\end{align*}
% END RESIDUAL 59
% BEGIN RESIDUAL 60
\begin{align*}
R_{60} ={}& -x_{9} x_{62} - x_{13} + x_{75} x_{76} + x_{77}
\end{align*}
% END RESIDUAL 60
% BEGIN RESIDUAL 61
\begin{align*}
R_{61} ={}& -1 + 8x_{2} + 4x_{2} x_{8} + 8x_{2} x_{9} - 2x_{2} x_{19} + x_{75}
\end{align*}
% END RESIDUAL 61
% BEGIN RESIDUAL 62
\begin{align*}
R_{62} ={}& -x_{8} x_{201} - x_{12} + x_{76}
\end{align*}
% END RESIDUAL 62
% BEGIN RESIDUAL 63
\begin{align*}
R_{63} ={}& x_{9} x_{201} - 2x_{19} x_{62} + 2x_{32} x_{62} + x_{67} + x_{78}
\end{align*}
% END RESIDUAL 63
% BEGIN RESIDUAL 64
\begin{align*}
R_{64} ={}& -x_{12} x_{63} + x_{13} x_{63} + x_{79}
\end{align*}
% END RESIDUAL 64
% BEGIN RESIDUAL 65
\begin{align*}
R_{65} ={}& -x_{64} x_{79} + x_{80}
\end{align*}
% END RESIDUAL 65
% BEGIN RESIDUAL 66
\begin{align*}
R_{66} ={}& -x_{65} x_{80} + x_{81}
\end{align*}
% END RESIDUAL 66
% BEGIN RESIDUAL 67
\begin{align*}
R_{67} ={}& -x_{66} x_{81} + x_{82}
\end{align*}
% END RESIDUAL 67
% BEGIN RESIDUAL 68
\begin{align*}
R_{68} ={}& -x_{67} x_{82} + x_{83}
\end{align*}
% END RESIDUAL 68
% BEGIN RESIDUAL 69
\begin{align*}
R_{69} ={}& -x_{68} x_{83} + x_{84}
\end{align*}
% END RESIDUAL 69
% BEGIN RESIDUAL 70
\begin{align*}
R_{70} ={}& -x_{69} x_{84} + x_{85}
\end{align*}
% END RESIDUAL 70
% BEGIN RESIDUAL 71
\begin{align*}
R_{71} ={}& -x_{70} x_{85} + x_{86}
\end{align*}
% END RESIDUAL 71
% BEGIN RESIDUAL 72
\begin{align*}
R_{72} ={}& -x_{77} x_{86} + x_{87}
\end{align*}
% END RESIDUAL 72
% BEGIN RESIDUAL 73
\begin{align*}
R_{73} ={}& -x_{78} x_{87} + x_{88}
\end{align*}
% END RESIDUAL 73
% BEGIN RESIDUAL 74
\begin{align*}
R_{74} ={}& -x_{88}^{2} + x_{89}
\end{align*}
% END RESIDUAL 74
% BEGIN RESIDUAL 75
\begin{align*}
R_{75} ={}& -x_{4} - 2x_{4} x_{16} + 2x_{16} + x_{90}
\end{align*}
% END RESIDUAL 75
% BEGIN RESIDUAL 76
\begin{align*}
R_{76} ={}& -x_{14} x_{90} + x_{91}
\end{align*}
% END RESIDUAL 76
% BEGIN RESIDUAL 77
\begin{align*}
R_{77} ={}& -2x_{1}^{2} + 2x_{1} x_{12} + 2x_{1} x_{13} + 2x_{1} x_{91} + 2x_{1} x_{92} - x_{12}^{2} \\
&{}  - 2x_{12} x_{91} - x_{13}^{2} - 2x_{13} x_{92} + x_{17} x_{93} - x_{73} x_{89} - x_{91}^{2} \\
&{}  - x_{92}^{2}
\end{align*}
% END RESIDUAL 77
% BEGIN RESIDUAL 78
\begin{align*}
R_{78} ={}& -x_{4} x_{15} + x_{15} - x_{15} x_{90} + x_{92}
\end{align*}
% END RESIDUAL 78
% BEGIN RESIDUAL 79
\begin{align*}
R_{79} ={}& -x_{94} x_{97} + x_{97} + x_{99}
\end{align*}
% END RESIDUAL 79
% BEGIN RESIDUAL 80
\begin{align*}
R_{80} ={}& -x_{3} x_{120} + x_{95} x_{193} - x_{98} - x_{99} x_{120}
\end{align*}
% END RESIDUAL 80
% BEGIN RESIDUAL 81
\begin{align*}
R_{81} ={}& x_{100} x_{193} - x_{101} x_{123} - x_{102}
\end{align*}
% END RESIDUAL 81
% BEGIN RESIDUAL 82
\begin{align*}
R_{82} ={}& x_{103} x_{193} - x_{104} x_{123} - x_{105}
\end{align*}
% END RESIDUAL 82
% BEGIN RESIDUAL 83
\begin{align*}
R_{83} ={}& x_{106} x_{193} - x_{107} x_{123} - x_{108}
\end{align*}
% END RESIDUAL 83
% BEGIN RESIDUAL 84
\begin{align*}
R_{84} ={}& x_{109} x_{193} - x_{110} x_{123} - x_{111}
\end{align*}
% END RESIDUAL 84
% BEGIN RESIDUAL 85
\begin{align*}
R_{85} ={}& x_{112} x_{194} - x_{114} x_{123} - x_{115}
\end{align*}
% END RESIDUAL 85
% BEGIN RESIDUAL 86
\begin{align*}
R_{86} ={}& x_{116} x_{194} - x_{117} x_{123} - x_{118}
\end{align*}
% END RESIDUAL 86
% BEGIN RESIDUAL 87
\begin{align*}
R_{87} ={}& -x_{17} x_{120} + x_{123}
\end{align*}
% END RESIDUAL 87
% BEGIN RESIDUAL 88
\begin{align*}
R_{88} ={}& x_{119} - x_{121} x_{123} - x_{122}
\end{align*}
% END RESIDUAL 88
% BEGIN RESIDUAL 89
\begin{align*}
R_{89} ={}& -x_{123} x_{125} + x_{124} - x_{126}
\end{align*}
% END RESIDUAL 89
% BEGIN RESIDUAL 90
\begin{align*}
R_{90} ={}& -x_{123} x_{128} + x_{127} x_{193} - x_{129}
\end{align*}
% END RESIDUAL 90
% BEGIN RESIDUAL 91
\begin{align*}
R_{91} ={}& x_{130} - 2x_{130} x_{132} + x_{133}
\end{align*}
% END RESIDUAL 91
% BEGIN RESIDUAL 92
\begin{align*}
R_{92} ={}& -x_{130} x_{133} - x_{133} + x_{134}
\end{align*}
% END RESIDUAL 92
% BEGIN RESIDUAL 93
\begin{align*}
R_{93} ={}& -1 + x_{131}^{2} + x_{133}^{2} - x_{134}^{2}
\end{align*}
% END RESIDUAL 93
% BEGIN RESIDUAL 94
\begin{align*}
R_{94} ={}& 1 - x_{131} x_{136} - x_{135} + x_{137}
\end{align*}
% END RESIDUAL 94
% BEGIN RESIDUAL 95
\begin{align*}
R_{95} ={}& -1 + 2x_{135} - x_{135}^{2} + 2x_{135} x_{137} + x_{136}^{2}
\end{align*}
% END RESIDUAL 95
% BEGIN RESIDUAL 96
\begin{align*}
R_{96} ={}& 1 + x_{3}^{2} - 2x_{3} x_{131} + x_{138}
\end{align*}
% END RESIDUAL 96
% BEGIN RESIDUAL 97
\begin{align*}
R_{97} ={}& -x_{3} x_{170} + x_{17} + x_{138} x_{139} - x_{169}
\end{align*}
% END RESIDUAL 97
% BEGIN RESIDUAL 98
\begin{align*}
R_{98} ={}& -5 - 8x_{2} - 4x_{2}^{2} - 16x_{2} x_{16} + 4x_{2} x_{131} - 16x_{16} \\
&{}  - 16x_{16}^{2} + 8x_{16} x_{131} + 4x_{131} - x_{140}
\end{align*}
% END RESIDUAL 98
% BEGIN RESIDUAL 99
\begin{align*}
R_{99} ={}& -1 - 2x_{2} x_{170} - 4x_{16} x_{170} + x_{94} + x_{140} x_{141} - x_{169} \\
&{}  - 2x_{170}
\end{align*}
% END RESIDUAL 99
% BEGIN RESIDUAL 100
\begin{align*}
R_{100} ={}& -x_{131} x_{143} - x_{142} + x_{144}
\end{align*}
% END RESIDUAL 100
% BEGIN RESIDUAL 101
\begin{align*}
R_{101} ={}& -1 - x_{142}^{2} + 2x_{142} x_{144} + x_{143}^{2}
\end{align*}
% END RESIDUAL 101
% BEGIN RESIDUAL 102
\begin{align*}
R_{102} ={}& -x_{131} x_{146} - x_{145} + x_{147}
\end{align*}
% END RESIDUAL 102
% BEGIN RESIDUAL 103
\begin{align*}
R_{103} ={}& -1 - x_{145}^{2} + 2x_{145} x_{147} + x_{146}^{2}
\end{align*}
% END RESIDUAL 103
% BEGIN RESIDUAL 104
\begin{align*}
R_{104} ={}& -x_{131} x_{149} - x_{148} + x_{150}
\end{align*}
% END RESIDUAL 104
% BEGIN RESIDUAL 105
\begin{align*}
R_{105} ={}& -1 - x_{148}^{2} + 2x_{148} x_{150} + x_{149}^{2}
\end{align*}
% END RESIDUAL 105
% BEGIN RESIDUAL 106
\begin{align*}
R_{106} ={}& -x_{131} x_{152} - x_{151} + x_{153}
\end{align*}
% END RESIDUAL 106
% BEGIN RESIDUAL 107
\begin{align*}
R_{107} ={}& -1 - x_{151}^{2} + 2x_{151} x_{153} + x_{152}^{2}
\end{align*}
% END RESIDUAL 107
% BEGIN RESIDUAL 108
\begin{align*}
R_{108} ={}& -x_{131} x_{155} - x_{154} + x_{156}
\end{align*}
% END RESIDUAL 108
% BEGIN RESIDUAL 109
\begin{align*}
R_{109} ={}& -1 - x_{154}^{2} + 2x_{154} x_{156} + x_{155}^{2}
\end{align*}
% END RESIDUAL 109
% BEGIN RESIDUAL 110
\begin{align*}
R_{110} ={}& -x_{131} x_{158} - x_{157} + x_{159}
\end{align*}
% END RESIDUAL 110
% BEGIN RESIDUAL 111
\begin{align*}
R_{111} ={}& -1 - x_{157}^{2} + 2x_{157} x_{159} + x_{158}^{2}
\end{align*}
% END RESIDUAL 111
% BEGIN RESIDUAL 112
\begin{align*}
R_{112} ={}& 1 - x_{131} x_{161} - x_{160} + x_{162}
\end{align*}
% END RESIDUAL 112
% BEGIN RESIDUAL 113
\begin{align*}
R_{113} ={}& -1 + 2x_{160} - x_{160}^{2} + 2x_{160} x_{162} + x_{161}^{2}
\end{align*}
% END RESIDUAL 113
% BEGIN RESIDUAL 114
\begin{align*}
R_{114} ={}& -x_{131} x_{164} - x_{163} + x_{165}
\end{align*}
% END RESIDUAL 114
% BEGIN RESIDUAL 115
\begin{align*}
R_{115} ={}& -1 - x_{163}^{2} + 2x_{163} x_{165} + x_{164}^{2}
\end{align*}
% END RESIDUAL 115
% BEGIN RESIDUAL 116
\begin{align*}
R_{116} ={}& -x_{131} x_{167} - x_{166} + x_{168}
\end{align*}
% END RESIDUAL 116
% BEGIN RESIDUAL 117
\begin{align*}
R_{117} ={}& -1 - x_{166}^{2} + 2x_{166} x_{168} + x_{167}^{2}
\end{align*}
% END RESIDUAL 117
% BEGIN RESIDUAL 118
\begin{align*}
R_{118} ={}& 1 - x_{131} x_{170} - x_{169} + x_{171}
\end{align*}
% END RESIDUAL 118
% BEGIN RESIDUAL 119
\begin{align*}
R_{119} ={}& -1 + 2x_{169} - x_{169}^{2} + 2x_{169} x_{171} + x_{170}^{2}
\end{align*}
% END RESIDUAL 119
% BEGIN RESIDUAL 120
\begin{align*}
R_{120} ={}& -x_{131} x_{173} - x_{172} + x_{174}
\end{align*}
% END RESIDUAL 120
% BEGIN RESIDUAL 121
\begin{align*}
R_{121} ={}& -1 - x_{172}^{2} + 2x_{172} x_{174} + x_{173}^{2}
\end{align*}
% END RESIDUAL 121
% BEGIN RESIDUAL 122
\begin{align*}
R_{122} ={}& -x_{131} x_{176} - x_{175} + x_{177}
\end{align*}
% END RESIDUAL 122
% BEGIN RESIDUAL 123
\begin{align*}
R_{123} ={}& -1 - x_{175}^{2} + 2x_{175} x_{177} + x_{176}^{2}
\end{align*}
% END RESIDUAL 123
% BEGIN RESIDUAL 124
\begin{align*}
R_{124} ={}& -8 - 2x_{0} - 4x_{0}^{2} - 2x_{0} x_{2} - x_{1}^{2} - 8x_{2} \\
&{}  - 6x_{2}^{2} - 16x_{2} x_{16} - x_{3}^{2} - x_{4}^{2} - x_{8}^{2} - x_{9}^{2} \\
&{}  - x_{10}^{2} - x_{11}^{2} - x_{12}^{2} - x_{13}^{2} - x_{14}^{2} - x_{15}^{2} \\
&{}  - 16x_{16} - 20x_{16}^{2} - x_{17}^{2} + 2x_{94} - 2x_{94}^{2} - x_{95}^{2} \\
&{}  - x_{96}^{2} - x_{100}^{2} - x_{103}^{2} - x_{106}^{2} - x_{109}^{2} - x_{112}^{2} \\
&{}  - x_{113}^{2} - x_{116}^{2} - x_{119}^{2} - x_{120}^{2} - x_{124}^{2} - x_{127}^{2} \\
&{}  + x_{130}
\end{align*}
% END RESIDUAL 124
% BEGIN RESIDUAL 125
\begin{align*}
R_{125} ={}& -16 + x_{0} - 4x_{178}^{2}
\end{align*}
% END RESIDUAL 125
% BEGIN RESIDUAL 126
\begin{align*}
R_{126} ={}& -x_{131} x_{180} + x_{189}
\end{align*}
% END RESIDUAL 126
% BEGIN RESIDUAL 127
\begin{align*}
R_{127} ={}& -1 + x_{179}^{2} + x_{180}^{2} - x_{189}^{2}
\end{align*}
% END RESIDUAL 127
% BEGIN RESIDUAL 128
\begin{align*}
R_{128} ={}& -x_{182} x_{183} + x_{190}
\end{align*}
% END RESIDUAL 128
% BEGIN RESIDUAL 129
\begin{align*}
R_{129} ={}& -1 + x_{181}^{2} + x_{182}^{2} - x_{190}^{2}
\end{align*}
% END RESIDUAL 129
% BEGIN RESIDUAL 130
\begin{align*}
R_{130} ={}& -x_{176}^{2} + x_{191}
\end{align*}
% END RESIDUAL 130
% BEGIN RESIDUAL 131
\begin{align*}
R_{131} ={}& x_{180} - 2x_{185} x_{191} - x_{191}
\end{align*}
% END RESIDUAL 131
% BEGIN RESIDUAL 132
\begin{align*}
R_{132} ={}& -1 - 4x_{180} - 8x_{180} x_{184} + x_{183}
\end{align*}
% END RESIDUAL 132
% BEGIN RESIDUAL 133
\begin{align*}
R_{133} ={}& -x_{131} - x_{179} x_{186} + x_{183}
\end{align*}
% END RESIDUAL 133
% BEGIN RESIDUAL 134
\begin{align*}
R_{134} ={}& -x_{177} - x_{179} x_{187} + x_{181}
\end{align*}
% END RESIDUAL 134
% BEGIN RESIDUAL 135
\begin{align*}
R_{135} ={}& -2x_{16} - 4x_{176} x_{188} + x_{182}
\end{align*}
% END RESIDUAL 135
% BEGIN RESIDUAL 136
\begin{align*}
R_{136} ={}& 9 + 2x_{0}^{2} + 2x_{0} x_{2} - 2x_{0} x_{136} - 2x_{0} x_{143} + x_{1}^{2} \\
&{}  - 2x_{1} x_{170} + x_{2}^{2} - 2x_{2} x_{143} + x_{4}^{2} - 2x_{4} x_{161} + x_{8}^{2} \\
&{}  - 2x_{8} x_{146} + x_{9}^{2} - 2x_{9} x_{149} + x_{10}^{2} - 2x_{10} x_{152} + x_{11}^{2} \\
&{}  - 2x_{11} x_{155} + x_{12}^{2} - 2x_{12} x_{158} + x_{13}^{2} - 2x_{13} x_{164} + x_{14}^{2} \\
&{}  - 2x_{14} x_{167} + x_{15}^{2} - 2x_{15} x_{173} - 4x_{16} x_{176} - x_{131} x_{192} + x_{137}^{2} \\
&{}  - 2x_{144} + x_{144}^{2} - 2x_{147} + x_{147}^{2} - 2x_{150} + x_{150}^{2} \\
&{}  - 2x_{153} + x_{153}^{2} - 2x_{156} + x_{156}^{2} - 2x_{159} + x_{159}^{2} \\
&{}  + x_{162}^{2} - 2x_{165} + x_{165}^{2} - 2x_{168} + x_{168}^{2} + x_{171}^{2} \\
&{}  - 2x_{174} + x_{174}^{2} + 2x_{176}^{2} - 2x_{177} + x_{177}^{2} + x_{192} \\
&{}  - x_{195}^{2} - x_{196}^{2} - x_{197} - x_{197}^{2}
\end{align*}
% END RESIDUAL 136
% BEGIN RESIDUAL 137
\begin{align*}
R_{137} ={}& -x_{96} x_{193} + x_{120}
\end{align*}
% END RESIDUAL 137
% BEGIN RESIDUAL 138
\begin{align*}
R_{138} ={}& -x_{113} x_{194} + x_{120}
\end{align*}
% END RESIDUAL 138
% BEGIN RESIDUAL 139
\begin{align*}
R_{139} ={}& -1 - 4x_{16}^{2} - x_{136}^{2} - x_{143}^{2} - x_{146}^{2} - x_{149}^{2} \\
&{}  - x_{152}^{2} - x_{155}^{2} - x_{158}^{2} - x_{161}^{2} - x_{164}^{2} - x_{167}^{2} \\
&{}  - x_{170}^{2} - x_{173}^{2} + x_{176}^{2} - x_{195}^{2} - x_{196}^{2} - x_{197} \\
&{}  - x_{197}^{2}
\end{align*}
% END RESIDUAL 139
% BEGIN RESIDUAL 140
\begin{align*}
R_{140} ={}& -3 - 4x_{98}^{2} - 4x_{102}^{2} - 4x_{105}^{2} - 4x_{108}^{2} - 4x_{111}^{2} \\
&{}  - 4x_{115}^{2} - 4x_{118}^{2} + x_{120}^{2} - 4x_{122}^{2} - 4x_{126}^{2} - 4x_{129}^{2} \\
&{}  - x_{198}^{2} - x_{199}^{2} + 3x_{200} - x_{200}^{2}
\end{align*}
% END RESIDUAL 140
% BEGIN RESIDUAL 141
\begin{align*}
R_{141} ={}& -x_{2} x_{12} + x_{201}
\end{align*}
% END RESIDUAL 141
% BEGIN RESIDUAL 142
\begin{align*}
R_{142} ={}& x_{94} x_{203} - 2x_{95} x_{143} - 2x_{100} x_{146} - 2x_{103} x_{149} - 2x_{106} x_{152} - 2x_{109} x_{155} \\
&{}  - 2x_{112} x_{158} - 2x_{116} x_{164} - 2x_{119} x_{167} - 2x_{124} x_{173} - 2x_{127} x_{176} - 2x_{131} x_{203} \\
&{}  + 2x_{143} x_{144} + 2x_{146} x_{147} + 2x_{149} x_{150} + 2x_{152} x_{153} + 2x_{155} x_{156} + 2x_{158} x_{159} \\
&{}  + 2x_{164} x_{165} + 2x_{167} x_{168} + 2x_{173} x_{174} + 2x_{176} x_{177} + x_{202}
\end{align*}
% END RESIDUAL 142
% BEGIN RESIDUAL 143
\begin{align*}
R_{143} ={}& -10 - x_{94} x_{202} + x_{95}^{2} - 2x_{95} x_{142} + x_{100}^{2} - 2x_{100} x_{145} \\
&{}  + x_{103}^{2} - 2x_{103} x_{148} + x_{106}^{2} - 2x_{106} x_{151} + x_{109}^{2} - 2x_{109} x_{154} \\
&{}  + x_{112}^{2} - 2x_{112} x_{157} + x_{116}^{2} - 2x_{116} x_{163} + x_{119}^{2} - 2x_{119} x_{166} \\
&{}  + x_{124}^{2} - 2x_{124} x_{172} + x_{127}^{2} - 2x_{127} x_{175} + 2x_{142} x_{144} + 2x_{145} x_{147} \\
&{}  + 2x_{148} x_{150} + 2x_{151} x_{153} + 2x_{154} x_{156} + 2x_{157} x_{159} + 2x_{163} x_{165} + 2x_{166} x_{168} \\
&{}  + 2x_{172} x_{174} + 2x_{175} x_{177} + x_{203}
\end{align*}
% END RESIDUAL 143
% BEGIN RESIDUAL 144
\begin{align*}
R_{144} ={}& x_{94} x_{205} - 2x_{96} x_{136} - 2x_{113} x_{161} - 2x_{120} x_{170} - 2x_{131} x_{205} + 2x_{136} \\
&{}  + 2x_{136} x_{137} + 2x_{161} + 2x_{161} x_{162} + 2x_{170} + 2x_{170} x_{171} + x_{204} \\
&{}  - x_{205}
\end{align*}
% END RESIDUAL 144
% BEGIN RESIDUAL 145
\begin{align*}
R_{145} ={}& -2 - x_{94} x_{204} + x_{96}^{2} - 2x_{96} x_{135} + x_{113}^{2} - 2x_{113} x_{160} \\
&{}  + x_{120}^{2} - 2x_{120} x_{169} + 2x_{135} + 2x_{135} x_{137} + 2x_{160} + 2x_{160} x_{162} \\
&{}  + x_{169}^{2} - x_{170}^{2} + x_{204} + x_{205}
\end{align*}
% END RESIDUAL 145
% BEGIN RESIDUAL 146
\begin{align*}
R_{146} ={}& -1 - x_{2} - 2x_{2} x_{206} + x_{4}
\end{align*}
% END RESIDUAL 146
% BEGIN RESIDUAL 147
\begin{align*}
R_{147} ={}& -x_{0} x_{2} + x_{208}
\end{align*}
% END RESIDUAL 147
% BEGIN RESIDUAL 148
\begin{align*}
R_{148} ={}& -1 - 8x_{2} x_{208} + x_{207}^{2} - 16x_{208}^{2}
\end{align*}
% END RESIDUAL 148
% BEGIN RESIDUAL 149
\begin{align*}
R_{149} ={}& -524288x_{0} - 1048576x_{0} x_{209} + x_{2}
\end{align*}
% END RESIDUAL 149
}

\endgroup
\end{document}
